% Prose for this counterexample.  Every region below is delimited by an
% environment that tools/build.py extracts; machine facts (ids, uid, dates,
% urls, enums) live in case.json.  See counterexamples/AGENTS.md.

\cxtitle{Borcea--Br\"and\'en AIM problems}

% Four problems from one list, one section.  Each pair below has its own
% witness, its own certificate module (verify_pot.py, verify_pencil.py) and its
% own finding history, so each opens with \cxsubtitle and its credit line.

\cxsubtitle{Problems 35 and 38: mixed-determinant representation and Permanent-On-Top}

\begin{cxcredits}
\posedby{Borcea and Br\"and\'en \citeyearpar{AIMStableProblems}}
\foundby{GPT-5 (Pro)}{2026-01-10}
\formalizedby{S.~Sra}
\auditedby{S.~Sra}
\contributedby{S.~Sra}
\end{cxcredits}

\begin{cxcontext}
A polynomial $f\in\R[z_1,\ldots,z_n]$ is \emph{real stable} if $f(z_1,\ldots,z_n)\ne0$ whenever every $z_i$ has strictly positive imaginary part.  A homogeneous symmetric polynomial of degree $d$ expands uniquely in the basis of Schur polynomials as $p=\sum_{\lambda\vdash d}a_\lambda s_\lambda$; we say $p$ is \emph{Schur-positive} if every $a_\lambda\ge0$.  For a partition $\lambda\vdash d$ we write $f^\lambda$ for the number of standard Young tableaux of shape $\lambda$, $\lambda'$ for the conjugate partition, and, for a $d\times d$ matrix $A=(a_{ij})$,
\[
\Imm_\lambda(A) :=\sum_{\sigma\in S_d}\chi^\lambda(\sigma)\prod_{i=1}^da_{i\sigma(i)},
\]
for the immanant induced by the irreducible character $\chi^\lambda$ of $S_d$. The two extreme cases are $\Imm_{(1^d)}=\det$ and $\Imm_{(d)}=\per$.  Next, for $n\times n$ matrices $A_1,\ldots,A_m$ the \emph{mixed-determinant} is
\[
\eta(A_1,\ldots,A_m) :=\sum_{(S_1,\ldots,S_m)}\det(A_1[S_1])\cdots\det(A_m[S_m]),
\]
where the sum runs over all ordered partitions of $\{1,\ldots,n\}$ into disjoint sets $S_1,\ldots,S_m$, and $A[S]$ denotes the principal submatrix of $A$ with rows and columns in $S$.  When $A\succeq0$ the polynomial $\eta(z_1A,\ldots,z_nA)$ is symmetric, homogeneous, real stable, and has the (Schur-positive) representation
\begin{equation*}
  \eta(z_1A,\ldots,z_nA)=\sum_{\lambda\vdash n}\Imm_{\lambda'}(A)s_\lambda(z_1,\ldots,z_n).
\end{equation*}
The two problems below ask how far those properties go the other way.
\end{cxcontext}

% ---------------- result: aim-problem-35 ----------------

\begin{cxsource}{aim-problem-35}
Borcea--Br\"and\'en, AIM problem list \citeyearpar{AIMStableProblems}, Problem 35
\end{cxsource}

\begin{cxstatement}{aim-problem-35}
Characterize all symmetric, real stable and homogeneous polynomial of degree $d$ in $n$ variables.  Is it in fact so that if $p$ is such a polynomial, then there is a $d\times d$ positive semidefinite $d\times d$ matrix $A$ such that
\[p(z_1,\ldots,z_n)=\eta(z_1A,\ldots,z_nA)?\]
This is not obvious even for $n=2$.
\end{cxstatement}

\begin{cxsummary}{aim-problem-35}
Borcea--Br\"and\'en Problem 35: common-matrix mixed-determinant representation \citeyearpar{AIMStableProblems}
\end{cxsummary}

\begin{cxcertificate}{aim-problem-35}
Stable degree-five polynomial plus size-five immanant POT obstruction.
\end{cxcertificate}

% ---------------- result: aim-problem-38 ----------------

\begin{cxsource}{aim-problem-38}
Borcea--Br\"and\'en, AIM problem list \citeyearpar{AIMStableProblems}, Problem 38
\end{cxsource}

\begin{cxstatement}{aim-problem-38}
Lieb's Permanent-On-Top (POT) Conjecture for the symmetric group on $d$ elements [J,L] asserts that if $A$ is a $d\times d$ positive semi-definite matrix then $\Imm_\lambda(A)\le f^\lambda\per(A)$ for all $\lambda\vdash d$.  An affirmative answer to the following problem would in particular imply the validity of the POT Conjecture.

\emph{Problem 38.}  Let $p$ be as in Problem 36.  Is $a_\lambda\le f^\lambda a_{(1^d)}$ for all $\lambda\vdash d$?
\end{cxstatement}

\begin{cxsummary}{aim-problem-38}
Borcea--Br\"and\'en Problem 38: Stable Permanent-On-Top \citeyearpar{AIMStableProblems}
\end{cxsummary}

\begin{cxcertificate}{aim-problem-38}
Exact Schur expansion; four violating partitions.
\end{cxcertificate}

\begin{cxrefutation}
We answers to the above problems are negative. Let $S=x_1+\cdots+x_5$ and define the polynomial
\begin{equation}\label{eq:stable-p}
p(x_1,\ldots,x_5)=\prod_{i=1}^5\Bigl(x_i+5\nlsum_{j\ne i}x_j\Bigr)=\prod_{i=1}^5(5S-4x_i).
\end{equation}
Every factor has strictly positive coefficients, hence has positive imaginary part whenever all variables do; therefore $p$ is real stable. Moreover, $p$ is symmetric, homogeneous of degree five in five variables, and has strictly positive monomial coefficients. % (it will also turn out to be Schur positive).

\begin{theorem}[\statusfalse: Stable Permanent-On-Top, AIM Problem 38]\label{thm:pot}
The polynomial \eqref{eq:stable-p} satisfies all hypotheses of Problem~\ref{ctx:aim-problem-38} but violates $a_\lambda\le f^\lambda a_{(1^5)}$.
\end{theorem}
\begin{proof}
A short explicit calculation shows that $p$ has the Schur expansion
\begin{align*}
p={}&625s_{(5)}+4500s_{(4,1)}+7125s_{(3,2)}+8150s_{(3,1,1)}+7525s_{(2,2,1)}+6580s_{(2,1,1,1)}+1281s_{(1^5)}.
\end{align*}
For $\lambda=(3,2)$, the hook-length formula gives $f^{(3,2)}=5$, whereas
\[
7125>5\cdot1281=6405.
\]
The same upper endpoint fails for $(3,1,1)$, $(2,2,1)$, and $(2,1,1,1)$.  The accompanying source package reconstructs the monomial expansion and inverts the Kostka matrix over $\mathbb Z$.
\end{proof}

The polynomial also has a diagonal positive-definite mixed-determinant representation.  For $j=1,\ldots,5$, let $D_j$ denote $5I$ except for the $j$-th entry being just 1. 
% \[
% D_j=\diag(5,\ldots,5,\underset{j}{1},5,\ldots,5).
% \]
Then the multivariate Johnson polynomial satisfies
\[
\eta(x_1D_1,\ldots,x_5D_5)=p(x_1,\ldots,x_5).
\]

\begin{theorem}[\statusfalse: common-matrix representation, AIM Problem 35]\label{thm:common-A}
There is no $5\times5$ Hermitian positive semidefinite matrix $A$ such that $p(x_1,\ldots,x_5)=\eta(x_1A,\ldots,x_5A)$.
\end{theorem}
\begin{proof}
Since $\eta(x_1A,\ldots,x_5A)$ is Schur positive with immanants as coefficients, the representation $p(x)=\eta(x)$ would identify the Schur coefficients (up to normalization). Since Lieb's permanental dominance inequality is known to hold in size five (indeed in size $\le 13$~\cite{Wanless2022} it is known), it would force
$a_\lambda\le f^\lambda a_{(1^5)}$ for every $\lambda\vdash5$. But this contradictions Theorem~\ref{thm:pot}, hence there cannot be any PSD $A$ for which $p(x)=\eta(x_1A,\ldots,x_5A)$.
\end{proof}

\begin{remark}[Scope]
The example is Schur-positive, and its lower endpoint inequalities hold, so it refutes neither Problem~\ref{ctx:aim-problem-36} nor Problem~\ref{ctx:aim-problem-37}; those require the positive-definite pencils of Theorems~\ref{thm:aim36} and~\ref{thm:aim37} below.  The full independent audit is included as \path{counterexamples/aim-problems/artifacts/}.
\end{remark}
\end{cxrefutation}

% ======================= Problems 36 and 37 =======================

\cxsubtitle{Problem 36: stable Schur positivity}

% Found in a later session than Problems 35 and 38; the other roles are
% inherited from the case block above.
\begin{cxcredits}[aim-problem-36]
\foundby{GPT-5.6 (Pro)}{2026-07-30}
\end{cxcredits}

\begin{cxrefutation}
  We keep the notation of the previous section. %: real stability means nonvanishing when every variable lies in the open upper half-plane; $a_\lambda$ are the coefficients in the Schur expansion $p=\sum_{\lambda\vdash d}a_\lambda s_\lambda$ of a symmetric homogeneous $p$ of degree $d$; $f^\lambda$ counts standard Young tableaux of shape $\lambda$; and $\Imm_\lambda$ is the immanant of a $d\times d$ matrix attached to the irreducible character $\chi^\lambda$, with $\Imm_{(1^d)}=\det$ and $\Imm_{(d)}=\per$.
Problem~\ref{ctx:aim-problem-36} asks whether stability forces a Schur-positive representation. %Problems~\ref{ctx:aim-problem-38} and~\ref{ctx:aim-problem-37} ask for the two endpoint inequalities that would bracket the Schur coefficients between the `$\det$' and `$\per$' ends of the immanant scale.
\end{cxrefutation}

% ---------------- result: aim-problem-36 ----------------

\begin{cxsource}{aim-problem-36}
Borcea--Br\"and\'en, AIM problem list \citeyearpar{AIMStableProblems}, Problem 36
\end{cxsource}

\begin{cxstatement}{aim-problem-36}
Let $p$ be a symmetric, real stable and homogeneous polynomial of degree $d$ in $n$ variables with $d\le n$ and suppose that $p$ has at least one positive (and then all non-negative) coefficients in the monomial basis.  Is $p$ Schur-positive?  That is, when $p(z_1,\ldots,z_n)$ is expanded in the Schur-basis
\[p(z_1,\ldots,z_n)=\sum_{\lambda\vdash d}a_\lambda s_\lambda(z_1,\ldots,z_n),\]
is $a_\lambda\ge0$ for all $\lambda\vdash d$?
\end{cxstatement}

\begin{cxsummary}{aim-problem-36}
Borcea--Br\"and\'en Problem 36: stable Schur positivity \citeyearpar{AIMStableProblems}
\end{cxsummary}

\begin{cxcertificate}{aim-problem-36}
Positive-definite $5\times5$ det-pencil; exact Schur expansion with $[s_{(1^5)}]q=-2972$.
\end{cxcertificate}

\begin{cxrefutation}
The answer to Problem~\ref{ctx:aim-problem-36} is no already at $d=n=5$.  Let
\begingroup
\scriptsize
\[
J_1=
\begin{pmatrix}
38&19&19&2&36\\
19&38&2&19&36\\
19&2&38&19&36\\
2&19&19&38&36\\
36&36&36&36&72
\end{pmatrix},\quad
J_2=
\begin{pmatrix}
38&19&19&17&21\\
19&38&17&19&21\\
19&17&38&34&36\\
17&19&34&38&36\\
21&21&36&36&42
\end{pmatrix},\quad
J_3=
\begin{pmatrix}
38&34&19&17&36\\
34&38&17&19&36\\
19&17&38&19&21\\
17&19&19&38&21\\
36&36&21&21&42
\end{pmatrix},
\]
\[
J_4=
\begin{pmatrix}
38&4&19&2&21\\
4&8&2&4&6\\
19&2&38&4&21\\
2&4&4&8&6\\
21&6&21&6&42
\end{pmatrix},\qquad
J_5=
\begin{pmatrix}
8&4&4&2&6\\
4&38&2&19&21\\
4&2&8&4&6\\
2&19&4&38&21\\
6&21&6&21&42
\end{pmatrix},
\]
\endgroup
and define the determinantal polynomial
\begin{equation}\label{eq:aim36-q}
q(x_1,\ldots,x_5)=\frac1{648}\det\Bigl(\nlsum_{i=1}^5x_iJ_i\Bigr).
\end{equation}

\begin{theorem}[\statusfalse: stable Schur positivity, AIM Problem 36]\label{thm:aim36}
The polynomial \eqref{eq:aim36-q} is symmetric, homogeneous of degree five in five variables, real stable, and all $126$ of its ordinary monomial coefficients are strictly positive, yet it is \emph{\textbf{not}} Schur positive, since
\[
[s_{(1^5)}]\,q=-2972<0.
\]
\end{theorem}
\begin{proof}
The leading principal minors of $J_i$ are
\[
\begin{array}{c|rrrrr}
 &\Delta_1&\Delta_2&\Delta_3&\Delta_4&\Delta_5\\ \hline
J_1&38&1083&28728&202176&1119744\\
J_2&38&1083&28728&202176&1119744\\
J_3&38&288&8208&202176&1119744\\
J_4&38&288&8208&46656&1119744\\
J_5&8&288&1728&46656&1119744,
\end{array}
\]
all positive, so Sylvester's criterion gives $J_i\succ0$; hence $q$ is real stable, and homogeneity of degree 5 is immediate. %If $z_1,\ldots,z_5$ lie in the open upper half-plane and $(\sum_iz_iJ_i)v=0$ for some $v\ne0$, then 
% \[
% 0=\operatorname{Im}\left\langle v,\left(\textstyle\sum_iz_iJ_i\right)v\right\rangle=\sum_i\operatorname{Im}(z_i)\langle v,J_iv\rangle>0,
% \]
% a contradiction; hence $q$ is real stable, and homogeneity of degree five is immediate.
Exact expansion of the determinant in the monomial basis gives
\begin{equation*}
  \begin{split}
    q={}&1728m_{(5)}+19440m_{(4,1)}+66555m_{(3,2)}+140910m_{(3,1,1)}\\
    &+250440m_{(2,2,1)}+547405m_{(2,1,1,1)}+1182860m_{(1^5)}.
  \end{split}
\end{equation*}
Inverting the Kostka matrix over $\mathbb Z$ in the partition order $(5),(4,1),(3,2),(3,1,1),(2,2,1),(2,1,1,1),(1^5)$ translates $q$ into the Schur basis, yielding
\begin{align*}
q={}&1728s_{(5)}+17712s_{(4,1)}+47115s_{(3,2)}+56643s_{(3,1,1)}+62415s_{(2,2,1)}+56437s_{(2,1,1,1)}-2972s_{(1^5)}.
\end{align*}
The last coefficient alone also follows from the single integer identity
\begin{align*}
[s_{(1^5)}]q={}&1728-2(19440)-2(66555)+3(140910)+3(250440)-4(547405)+1182860=-2972.
\end{align*}
The accompanying source package recomputes the Schur expansion in a second, independent way, and every step uses rational arithmetic only.
\end{proof}

\begin{remark}[Origin of the pencil]
The matrices are not a numerical accident.  Let $V=\ker B\cong S^{(3,2)}$, where $B$ maps the permutation module on $2$-subsets of $[5]$ to $\mathbb R^5$ by $e_{ab}\mapsto e_a+e_b$, let $P_i$ be the rank-three projection attached to the point stabilizer of $i$, and let $G$ be the Gram matrix of the chosen basis of $V$.  Then $J_i=2G(I+5P_i)$, and the whole one-parameter family $\det(\sum_ix_i(I+tP_i))$ has a closed Schur expansion whose $s_{(1^5)}$ coefficient turns negative for $t$ large; $t=5$ clears denominators.  \emph{Appendix~\ref{sec:aim36pencil} develops this construction in full.}
\end{remark}

\end{cxrefutation}

% ---------------- result: aim-problem-37 ----------------

\cxsubtitle{Problem 37: Schur-style inequality for stable polynomials}
% This result's witness has a different history from the case's: the pencil is
% the same GPT construction, but the Schur-positive member the lower endpoint
% actually needs was specialized out of it later.  Both are restated, because a
% per-result block replaces the inherited \foundby rather than appending to it.
\begin{cxcredits}[aim-problem-37]
\foundby{GPT-5.6 (Pro)}{2026-07-30}
\foundby{bugfixed by Opus 5}{2026-08-10}
\end{cxcredits}

Next, we close off the lower endpoint analog of Problem~\ref{ctx:aim-problem-38} (which concerned $a_\lambda \le f^\lambda a_{(1^d)}$).
\begin{cxsource}{aim-problem-37}
Borcea--Br\"and\'en, AIM problem list \citeyearpar{AIMStableProblems}, Problem 37
\end{cxsource}

\begin{cxstatement}{aim-problem-37}
If $A$ is a $d\times d$ matrix then $\det(A)=\Imm_{(1^d)}(A)$ and $\per(A)=\Imm_{(d)}(A)$, where $(1^d)=(1,\ldots,1)$ and $(d)=(1^d)'=(d,0,\ldots,0)$.  Schur's inequality [J,L] asserts that if $A$ is a $d\times d$ positive semi-definite matrix then $\Imm_\lambda(A)\ge f^\lambda\det(A)$ for all $\lambda\vdash d$, where $f^\lambda$ is the number of standard Young tableaux of shape $\lambda$.  A natural real stable extension of this result would be as follows.

\emph{Problem 37.}  Let $p$ be as in Problem 36.  Is $a_\lambda\ge f^\lambda a_{(d)}$ for all $\lambda\vdash d$?
\end{cxstatement}

\begin{cxsummary}{aim-problem-37}
Borcea--Br\"and\'en Problem 37: stable Schur inequality (lower endpoint) \citeyearpar{AIMStableProblems}
\end{cxsummary}

\begin{cxcertificate}{aim-problem-37}
Schur-positive member of the same family; $a_{(1^5)}=69<125=f^{(1^5)}a_{(5)}$.
\end{cxcertificate}


\begin{cxrefutation}
The lower endpoint needs a different member of the same family, because it needs a witness that is Schur \emph{positive}.  Let $G$ be the Gram matrix of the basis of $V$ that produced the $J_i$,
\[
G=\begin{pmatrix}
4&2&2&1&3\\
2&4&1&2&3\\
2&1&4&2&3\\
1&2&2&4&3\\
3&3&3&3&6
\end{pmatrix},
\]
put $\widetilde J_i:=(4J_i+2G)/5$, again a symmetric positive-definite \emph{integer} matrix, %, and in the notation of Appendix~\ref{sec:aim36pencil} equal to $2G(I+4P_i)$ ---
and define
\begin{equation}\label{eq:aim36-qtilde}
\widetilde q(x_1,\ldots,x_5)=\frac1{5184}\det\Bigl(\nlsum_{i=1}^5x_i\widetilde J_i\Bigr),
\end{equation}
which is $p_4$ in the notation of Appendix~\ref{sec:aim36pencil}.

\begin{theorem}[\statusfalse: lower endpoint, AIM Problem 37]\label{thm:aim37}
The polynomial \eqref{eq:aim36-qtilde} satisfies every hypothesis of Problem~\ref{ctx:aim-problem-37}, yet it violates the lower endpoint inequality $a_\lambda\ge f^\lambda a_{(d)}$: at $\lambda=(1^5)$, where $f^{(1^5)}=1$,
\[
a_{(1^5)}=69<125=a_{(5)}.
\]
\end{theorem}
\begin{proof}
The leading principal minors of the $\widetilde J_i$ are
\[
\begin{array}{c|rrrrr}
 &\Delta_1&\Delta_2&\Delta_3&\Delta_4&\Delta_5\\ \hline
\widetilde J_1&32&768&17280&118800&648000\\
\widetilde J_2&32&768&17280&118800&648000\\
\widetilde J_3&32&240&5760&118800&648000\\
\widetilde J_4&32&240&5760&32400&648000\\
\widetilde J_5&8&240&1440&32400&648000,
\end{array}
\]
all positive, so $\widetilde J_i\succ0$ and $\widetilde q$ is real stable, symmetric and homogeneous of degree five by the argument used for \eqref{eq:aim36-q}. It can be verified that $\widetilde{q}$ has the Schur positive representation
\begin{equation}
  \label{eq:1}
\widetilde{q}={}125s_{(5)}+1100s_{(4,1)}+2545s_{(3,2)}+3110s_{(3,1,1)}+3321s_{(2,2,1)}+2972s_{(2,1,1,1)}+69s_{(1^5)},
\end{equation}
which can be obtained by starting with the direct monomial basis expansion
\begin{align*}
\widetilde q={}&125m_{(5)}+1225m_{(4,1)}+3770m_{(3,2)}+7980m_{(3,1,1)}+13846m_{(2,2,1)}+30004m_{(2,1,1,1)}+64472m_{(1^5)},
\end{align*}
from which upon inverting the Kostka matrix in the same partition order yields~\eqref{eq:1}. Now $f^{(1^5)}=1$, so the lower endpoint at $\lambda=(1^5)$ would require $69\ge125$.  It is the only partition at which the inequality fails: the remaining six hold, the tightest being $\lambda=(4,1)$ with $1100\ge4\cdot125$.
\end{proof}

% \begin{remark}
% Problem~\ref{ctx:aim-problem-37} takes its hypotheses from Problem~\ref{ctx:aim-problem-36}, so $q$ satisfies them, and $[s_{(1^5)}]q=-2972<1728$ violates its conclusion; read literally, Theorem~\ref{thm:aim36} therefore refutes the lower endpoint as a corollary.  That corollary is empty.  Under those hypotheses $a_{(d)}$ coincides with the coefficient of $m_{(d)}$, hence with that of $x_1^d$, which is nonnegative --- so $a_\lambda\ge f^\lambda a_{(d)}$ forces $a_\lambda\ge0$, and Problem~\ref{ctx:aim-problem-37} as posed \emph{contains} Problem~\ref{ctx:aim-problem-36}.  Every counterexample to the latter is automatically one to the former, and says nothing further.  What the problem is asking, being the announced stable analogue of Schur's inequality $\Imm_\lambda(A)\ge f^\lambda\det(A)$ --- whose coefficients are immanants of a positive semidefinite matrix and so nonnegative by construction --- is whether the dominance survives among polynomials that \emph{are} Schur positive.  That is the reading Theorem~\ref{thm:aim37} answers, and because $\widetilde q$ is Schur positive it answers the literal reading too; nothing here turns on which was meant.
% \end{remark}

% \begin{remark}[Relation to Problem~\ref{ctx:aim-problem-38}]
% Both witnesses violate the upper endpoint as well: $q$ by the same sign obstruction, since at $\lambda=(5)$ it would require $1728\le-2972$, and $\widetilde q$ at six of the seven partitions, beginning with $125>1\cdot69$.  Neither refutation is carried here, because Theorem~\ref{thm:pot} already settles Problem~\ref{ctx:aim-problem-38} with a Schur-positive polynomial obtained independently of this pencil.
% \end{remark}
\end{cxrefutation}

%%% Local Variables:
%%% mode: LaTeX
%%% TeX-master: "../../tex/main"
%%% End:
